% Part: turing-machines
% Chapter: machines-computations
% Section: unary-numbers

\documentclass[../../../include/open-logic-section]{subfiles}

\begin{document}

\olfileid{tur}{mac}{una}
\olsection{సంఖ్యలకు ఏకాంక ప్రాతినిధ్యం}

\begin{explain}
ట్యూరింగ్ యంత్రాలు తమ టేపుపై వ్రాసిన సంకేతాల క్రమాలపై పనిచేస్తాయి. ట్యూరింగ్
యంత్రం ఉపయోగించే వర్ణమాలను బట్టి, ఈ సంకేతాల క్రమాలు వివిధ ఇన్‌పుట్‌లు,
అవుట్‌పుట్‌లకు ప్రాతినిధ్యం వహించగలవు. సహజంగానే, \emph{అంకగణిత}
ప్రమేయాలను, అంటే సహజ సంఖ్యల ప్రమేయాలను గణించే ట్యూరింగ్ యంత్రాలు ప్రత్యేక
ఆసక్తిని కలిగిస్తాయి. ధన పూర్ణ సంఖ్యలకు ప్రాతినిధ్యం వహించే ఒక సరళ పద్ధతి,
వాటిని ఒకే సంకేతం~$\TMstroke$ యొక్క క్రమాలుగా సంకేతీకరించడం. $n \in \Nat$
అయితే, $n = 0$ అయినప్పుడు $\TMstroke^n$ ఖాళీ క్రమం అనీ, లేకపోతే కచ్చితంగా
$n$ $\TMstroke$లతో కూడిన క్రమం అనీ ఉంచుదాం.
\end{explain}

\begin{defn}[గణన]
ట్యూరింగ్ యంత్రం~$M$ కింది ఇన్‌పుట్‌పై ఆగి
\[
\TMstroke^{n_1} \TMblank \TMstroke^{n_2} \TMblank \dots \TMblank \TMstroke^{n_k}
\]
$\TMstroke^{f(n_1, \dots, n_k)}$ను అవుట్‌పుట్‌గా ఇస్తే, అప్పుడే మాత్రమే
$M$ ప్రమేయం $f\colon \Nat^k \to \Nat$ను \emph{గణిస్తుంది}.
\end{defn}

\begin{prob}
  ట్యూరింగ్ యంత్రం~$M$ ప్రమేయం $f\colon \Nat^k \to \Nat^m$ను గణించడం
  ఎప్పుడో చెప్పే నిర్వచనాన్ని ఇవ్వండి.
\end{prob}

\begin{ex}\ollabel{ex:adder}
\emph{కూడిక:}
$f(n,m) = n + m$ అనే ప్రమేయాన్ని గణించే యంత్రాన్ని నిర్మిద్దాం. టేపుపై
పొడవులు $n$, $m$ అయిన రెండు $\TMstroke$ల దిమ్మలతో మొదలై, $n+m$
$\TMstroke$లతో కూడిన ఒకే దిమ్మతో ఆగే యంత్రం దీనికి కావాలి. ఇన్‌పుట్‌లోని
రెండు $\TMstroke$ల దిమ్మల మధ్య ఒక~$\TMblank$ ఉంటుంది; కాబట్టి $\TMblank$
ఉన్న గడిపై ఒక గీతను వ్రాసి, చివరి~$\TMstroke$ను చెరపడం ఒక పద్ధతి.
\begin{figure}\[
\begin{tikzpicture}[->,>=stealth',shorten >=1pt,auto,node distance=2.8cm,
                    semithick]
  \tikzstyle{every state}=[fill=none,draw=black,text=black]

  \node[initial,state]         (A)              {$q_0$};
  \node[state]         (B) [right of=A] {$q_1$};
  \node[state]         (C) [right of=B] {$q_2$};

  \path (A) edge node {\TMtrans{\TMblank}{\TMstroke}{\TMstay}} (B)
           edge [loop above] node {\TMtrans{\TMstroke}{\TMstroke}{\TMright}} (A)
        (B) edge [loop above] node {\TMtrans{\TMstroke}{\TMstroke}{\TMright}} (B)
            edge node {\TMtrans{\TMblank}{\TMblank}{\TMleft}} (C)
        (C) edge [loop above] node {\TMtrans{\TMstroke}{\TMblank}{\TMstay}} (C);
\end{tikzpicture}
\]\caption{$f(x,y) = x+y$ను గణించే యంత్రం}
\ollabel{fig:adder}
\end{figure}
\sourcecorrection{OLTETURMAC-005}{ఆంగ్ల మూలంలోని కూడిక యంత్రంలో స్థితి $q_0$ వద్ద $\TMstroke$ను చదివే బాణాన్ని ``loop''గా గీసినా దాని లక్ష్యాన్ని $q_1$గా ఇచ్చింది. అలా అయితే యంత్రం మొదటి ఇన్‌పుట్ దిమ్మలోని తొలి గీత తరువాతే రెండవ దశలోకి వెళ్లి కూడికను గణించదు. ఉద్దేశించిన స్వీయ-చక్రానికి అనుగుణంగా లక్ష్యాన్ని $q_0$గా సరిచేశాం.}
\end{ex}

\begin{prob}
ఇన్‌పుట్~$\tuple{3,2}$కు \olref[tur][mac][una]{ex:adder}లోని యంత్రపు
స్థితివిన్యాసాలను వరుసగా అనుసరించండి. యంత్రం $0+0$ను గణిస్తే ఏమి జరుగుతుంది?
\end{prob}

\begin{explain}
\olref[rep]{ex:doubler}లో $\TMstroke$ల క్రమాన్ని ఇన్‌పుట్‌గా తీసుకుని,
టేపుపై దానికి రెండింతల $\TMstroke$ల క్రమంతో ఆగే ట్యూరింగ్ యంత్రానికి---ద్విగుణక
యంత్రానికి---ఒక ఉదాహరణ ఇచ్చాం. అయితే, ద్విగుణీకరించిన $\TMstroke$ల దిమ్మకు
ఎడమవైపున అవుట్‌పుట్‌లో $\TMblank$లు ఉంటాయి కనుక, మీరు ఊహించి ఉండగల విధంగా
అది నిజానికి $f(x) = 2x$ ప్రమేయాన్ని గణించదు. దాన్ని సరిచేసే రెండు పద్ధతులను
వర్ణిస్తాం.
\end{explain}

\begin{ex}
\olref{fig:doubler-disc}లోని యంత్రం $f(x) = 2x$ ప్రమేయాన్ని గణిస్తుంది.
\olref[rep]{ex:doubler}లోని యంత్రం చేసినట్లు ఇన్‌పుట్‌ను చెరిపి, ఇన్‌పుట్‌లోని
ప్రతి $\TMstroke$కు అత్యంత కుడివైపున రెండు $\TMstroke$లను వ్రాయడానికి బదులుగా,
ఈ యంత్రం ఇన్‌పుట్‌లోని ప్రతి~$\TMstroke$కు కుడివైపున ఒకే~$\TMstroke$ను
చేర్చుతుంది. ఇన్‌పుట్ ఎక్కడ ముగుస్తుందో అది గుర్తుపెట్టుకోవాలి; అందుకే
ఇన్‌పుట్‌కూ చేర్చిన గీతలకూ మధ్య ఒక~$\TMblank$ను వదిలి, చివర్లో దాన్ని
ఒక~$\TMstroke$తో నింపుతుంది. అలాగే ఇన్‌పుట్‌లో మనం ఎక్కడ ఉన్నామో
``గుర్తుపెట్టుకోవాలి''; అందుకే ఇన్‌పుట్ దిమ్మలోని ఒక~$\TMstroke$ను తాత్కాలికంగా
ఒక~$\TMblank$తో మారుస్తాం.
  \begin{figure}
\[
\begin{tikzpicture}[->,>=stealth',shorten >=1pt,auto,node distance=2.8cm,
                    semithick]
  \tikzstyle{every state}=[fill=none,draw=black,text=black]

  \node[initial,state] (0)              {$q_0$};
  \node[state]         (1) [above of=0] {$q_1$};
  \node[state]         (2) [above of=1] {$q_2$};
  \node[state]         (3) [right of=2] {$q_3$};
  \node[state]         (4) [below of=3] {$q_4$};
  \node[state]         (5) [below of=4] {$q_5$};
  \node[state]         (6) [above right of=4] {$q_6$};
  \node[state]         (7) [below of=6] {$q_7$};
  \node[state]         (8) [below of=7] {$q_8$};

  \path (0) edge node {\TMtrans{\TMstroke}{\TMblank}{\TMright}} (1)
%    (0) edge node {\TMtrans{\TMblank}{\TMblank}{\TMstay}} (h)
    (1) edge [loop left] node {\TMtrans{\TMstroke}{\TMstroke}{\TMright}} (1)
      edge node {\TMtrans{\TMblank}{\TMblank}{\TMright}} (2)
    (2) edge [loop above] node {\TMtrans{\TMstroke}{\TMstroke}{\TMright}} (2)
        edge node {\TMtrans{\TMblank}{\TMstroke}{\TMleft}} (3)
    (3) edge [loop above] node {\TMtrans{\TMstroke}{\TMstroke}{\TMleft}} (3)
        edge node[left] {\TMtrans{\TMblank}{\TMblank}{\TMleft}} (4)
    (4) edge        node[left] {\TMtrans{\TMstroke}{\TMstroke}{\TMleft}} (5)
    (5) edge [loop below] node {\TMtrans{\TMstroke}{\TMstroke}{\TMleft}} (5)
        edge              node {\TMtrans{\TMblank}{\TMstroke}{\TMright}} (0)
    (4) edge      node[sloped] {\TMtrans{\TMblank}{\TMstroke}{\TMright}} (6)
    (6) edge              node {\TMtrans{\TMblank}{\TMstroke}{\TMright}} (7)
    (7) edge [loop right] node {\TMtrans{\TMstroke}{\TMstroke}{\TMright}} (7)
        edge              node {\TMtrans{\TMblank}{\TMblank}{\TMleft}} (8)
    (8) edge [loop right] node {\TMtrans{\TMstroke}{\TMblank}{\TMstay}} (8);
    \end{tikzpicture}
\]
\caption{$f(x) = 2x$ను గణించే యంత్రం}
\ollabel{fig:doubler-disc}
\end{figure}
\end{ex}

\begin{ex}\ollabel{ex:mover}
$f(x) = 2x$ను గణించడానికి రెండవ అవకాశం: మొదటి ద్విగుణక యంత్రాన్నే ఉంచి,
ద్విగుణీకరించిన గీతల దిమ్మను టేపులో అత్యంత ఎడమకు తరలించే స్థితులు,
నిర్దేశాలను చివర్లో చేర్చడం. \olref{fig:mover}లోని యంత్రం ఈ చివరి భాగాన్నే
చేస్తుంది: $\TMblank$ల దిమ్మ, దాని తరువాత $\TMstroke$ల దిమ్మ ఉన్న టేపుపై
(శీర్షం $\TMblank$ల దిమ్మలో ఎక్కడైనా ఉండగా) ప్రారంభిస్తే, అది $\TMstroke$లను
ఒక్కొక్కటిగా చెరిపి టేపు మొదట్లో వ్రాస్తుంది. పని ఎప్పుడు పూర్తయిందో
తెలుసుకోవడానికి, మొదట $\TMstroke$ల దిమ్మ చివరను ఒక~$\TMendtape$ సంకేతంతో
గుర్తిస్తుంది; చివర్లో ఆ సంకేతం చెరిగిపోతుంది. ఈ స్థితులను ద్విగుణక యంత్రానికి
చేర్చగలిగేలా వాటి సంఖ్యలను~$q_6$ వద్ద ప్రారంభించాం. అదనంగా కావలసిందల్లా
$\delta(q_0, \TMblank) = \tuple{q_6,\TMblank,\TMstay}$ అనే నిర్దేశం, అంటే
$\TMtrans{\TMblank}{\TMblank}{\TMstay}$ అనే గుర్తు గల బాణాన్ని~$q_0$ నుండి
~$q_6$కు చేర్చడం. (ఒక సూక్ష్మ సమస్య ఉంది: ఫలితంగా వచ్చే యంత్రం ఇన్‌పుట్~$x=0$కు
పనిచేయదు. దీనిని అభ్యాసంగా వదిలేస్తాం.)
\begin{figure}
  \[
  \begin{tikzpicture}[->,>=stealth',shorten >=1pt,auto,node distance=2.8cm,
                      semithick]
    \tikzstyle{every state}=[fill=none,draw=black,text=black]
    \node[initial,state] (6)              {$q_6$};
    \node[state]         (7) [right of=6] {$q_7$};
    \node[state]         (8) [right of=7] {$q_8$};
    \node[state]         (9) [below of=8] {$q_9$};
    \node[state]         (10) [left of=9] {$q_{10}$};
    \node[state]         (11) [left of=10]  {$q_{11}$};
    \node[state]         (12) [below of=11]  {$q_{12}$};
    \node[state]         (13) [right of=12] {$q_{13}$};
    \node[state]         (14) [right of=13] {$q_{14}$};
    \path
    (6)  edge node {\TMtrans{\TMblank}{\TMblank}{\TMright}} (7)
    (7)  edge [loop above] node {\TMtrans{\TMblank}{\TMblank}{\TMright}} (7)
         edge node {\TMtrans{\TMstroke}{\TMstroke}{\TMright}} (8)
    (8)  edge [loop above] node {\TMtrans{\TMstroke}{\TMstroke}{\TMright}} (8)
         edge node {\TMtrans{\TMblank}{\TMendtape}{\TMleft}} (9)
    (9)  edge [loop right] node {\TMtrans{\TMstroke}{\TMstroke}{\TMleft}} (9)
         edge node {\TMtrans{\TMblank}{\TMblank}{\TMright}} (10)
    (10) edge node[above] {\TMtrans{\TMstroke}{\TMblank}{\TMleft}} (11)
    (11) edge [loop above] node {\TMtrans{\TMblank}{\TMblank}{\TMleft}} (11)
         edge node[left] {\begin{tabular}{@{}l@{}}
          \TMtrans{\TMendtape}{\TMendtape}{\TMright}\\
          \TMtrans{\TMstroke}{\TMstroke}{\TMright}
         \end{tabular}} (12)
    (12) edge node[] {\TMtrans{\TMblank}{\TMstroke}{\TMright}} (13)
    (13) edge [loop below] node {\TMtrans{\TMblank}{\TMblank}{\TMright}} (13)
         edge node[sloped] {\TMtrans{\TMstroke}{\TMblank}{\TMleft}} (11)
    (13) edge node {\TMtrans{\TMendtape}{\TMblank}{\TMstay}} (14);
  \end{tikzpicture}
  \]
  \caption{$\TMstroke$ల దిమ్మను ఎడమకు తరలించడం}
  \ollabel{fig:mover}
\end{figure}
\end{ex}

\begin{prob}
\olref[tur][mac][una]{ex:mover}లో, \olref[tur][mac][una]{fig:doubler-disc}లోని
ద్విగుణక యంత్రం, \olref[tur][mac][una]{fig:mover}లోని తరలింపు యంత్రం కలిపిన
యంత్రాన్ని వర్ణించాం. ఈ సంయుక్త యంత్రాన్ని ఇన్‌పుట్~$x=0$పై, అంటే ఖాళీ టేపుపై
ప్రారంభిస్తే ఏమి జరుగుతుంది? ఈ సందర్భంలో యంత్రం అవుట్‌పుట్~$2x=0$తో ఆగేలా
దానిని ఎలా సరిచేస్తారు? (ఒక స్థితి, ఒక సంక్రమణను చేర్చి దీన్ని చేయగలగాలి.)
\end{prob}

\begin{prob}
\emph{తీసివేత:} $n > m$ అయ్యే పొడవులు $n$, $m$ గల, ఖాళీ కాని రెండు
గీతల సంకేతమాలలను ఇన్‌పుట్‌గా ఇచ్చినప్పుడు $f(n,m) = n - m$ ప్రమేయాన్ని
గణించే ట్యూరింగ్ యంత్రాన్ని రూపొందించండి.
\end{prob}

\begin{prob}
\emph{సమానత్వం:} కింది ప్రమేయాన్ని గణించే ట్యూరింగ్ యంత్రాన్ని రూపొందించండి:
\[
\fn{equality}(n,m) =
\begin{cases}
  \text{1} & \text{$n = m$ అయితే} \\
  \text{0} & \text{$n \neq m$ అయితే}
\end{cases}
\]
ఇక్కడ~$n$, $m \in \PosInt$.
\end{prob}

\begin{prob}
$x$, $y$లు టేపుపై ఒక $\TMblank$తో వేరుచేసిన $\TMstroke$ల సంకేతమాలలతో
సూచించిన ధన పూర్ణ సంఖ్యలైతే, $\min(x,y)$ ప్రమేయాన్ని గణించే ట్యూరింగ్
యంత్రాన్ని రూపొందించండి. యంత్రపు వర్ణమాలలో అదనపు సంకేతాలను ఉపయోగించవచ్చు.

$\min$ ప్రమేయం తన పరామితుల్లో అతి చిన్న విలువను ఎంచుకుంటుంది; అందువల్ల
$\min(3,5)=3$, $\min(20,16)=16$, $\min(4,4)=4$, తదితరాలు.
\end{prob}

\begin{defn}
  కింది షరతులు నెరవేరితే, అప్పుడే మాత్రమే ట్యూరింగ్ యంత్రం~$M$ పాక్షిక ప్రమేయం
  $f\colon \Nat^k \pto \Nat$ను గణిస్తుంది:
  \begin{enumerate}
    \item $f(n_1, \dots, n_k) = m$ అయితే, $M$ ఇన్‌పుట్
    $\TMstroke^{n_1}\concat\TMblank\concat
    \dots \concat\TMblank\concat\TMstroke^{n_k}$పై ఆగి
    $\TMstroke^{m}$ను అవుట్‌పుట్‌గా ఇస్తుంది.
    \item $f(n_1, \dots, n_k)$ నిర్వచితం కాకపోతే, $M$ అసలే ఆగదు లేదా
    ఒకే $\TMstroke$ల దిమ్మ కాని అవుట్‌పుట్‌తో ఆగుతుంది.
  \end{enumerate}
\end{defn}

\end{document}
